\section{Implementation and Reproducibility Details}
\label{app:implementation}
This working supplement distinguishes the implemented exploratory model from the matched training study still to be completed. Internal run identities and source hashes are maintained separately from the anonymous manuscript.

\paragraph{Architecture and training configuration.}
The current R4-AB implementation uses a randomly initialized ViT-Tiny visual encoder with patch size 14, a projector to width 192, and six shared transformer blocks with six heads and MLP width 768. Mode embeddings condition the shared transformer. The registered configuration specifies a batch size of 128, BF16 training, AdamW at learning rate $5\times10^{-5}$ with weight decay $10^{-3}$, and gradient clipping at 1.0. The evaluated weights are designated epoch 10, seed 3072. This identifies the evaluated artifacts; it does not assert that every checkpoint was trained from scratch with identical historical continuation. Resolved training history, data exposure, and learning-rate schedules must accompany the final matched study.

The implementation class also allocates parameters for unused C/D/E extensions. A/B-only training disables their objectives, but does not automatically remove their storage. Resource comparisons must distinguish total resident parameters, active parameters, and actual compute. No parameter-efficiency claim is derived from counting only the active heads.

\paragraph{Coupling details.}
At fit time, the predicted-action probability is $\min(e/10,0.5)$ with the effective epoch including any continuation offset. The inherited non-fitting loss path instead uses predicted actions with probability one. Validation loss is therefore not interchangeable with the fitting mixture. The current B timestep condition uses the action estimate's flow time, while latent supervision is downweighted at low flow times. Matched ablations must control the conditioning and weighting rather than silently change both with the input source.

\section{Controlled Training Comparisons}
\label{app:coupling}
\begin{table*}[t]
\centering\small
\begin{tabular}{p{0.20\textwidth}p{0.29\textwidth}p{0.42\textwidth}}
\toprule
Configuration & Changed connection & Question to be answered \\
\midrule
A-only & Remove dynamics objective & Does dynamics training improve the direct policy? \\
B-only & Remove action objective & Does action learning improve prediction or decisions? \\
Shared-GT & Shared A/B; recorded actions for B & What does shared multi-task training provide? \\
Shared-Pred-Detach & Mix estimated actions; stop their gradients & What is the additional effect of estimated-action exposure? \\
Full coupling & Retain gradients through estimated actions & Does the action-mediated gradient improve real decisions? \\
Shared encoder, two DiTs & Separate predictors; preserve input and gradient coupling & Is predictor sharing necessary under explicit resource budgets? \\
Independent A + B & Compose separately trained models at inference & Does the coupled system improve on modular composition? \\
\bottomrule
\end{tabular}
\caption{Required training contrasts, not completed performance results. Independent A+B reuses the single-task models. Shared-GT versus a single-task model changes both objectives and shared learning; it does not alone isolate parameter sharing.}
\label{tab:ablation_design}
\end{table*}

Each contrast should use matched data, seeds, goal definitions, checkpoint rules, and documented training-budget accounting. Fixed data exposure and fixed total compute answer different fairness questions. Independent models must encode observations with their own encoders instead of treating separately learned latent coordinates as interchangeable. Predictor comparisons use common candidate pools; policy comparisons hold the evaluation protocol fixed. At least three training seeds are planned for the core contrasts, with 200 independent evaluation episodes per task and seed as an initial target. These are design targets, not completed sample counts.

\section{Decision Diagnostics}
\label{app:diagnostics}
\paragraph{Fixed-pool ranking.}
The diagnostic pools use 50 restored initial states, six flow-step settings, and 256 candidates per setting. State restoration includes the registered simulator and random-state handling. Shared initial-noise identities support comparisons across flow steps. The reported aggregate is an equal mean across settings; settings from the same state remain together during bootstrap. Random success is estimated by the diagnostic's uniform candidate sampling; pool coverage means that at least one candidate succeeds. Physical-distance oracle and success-coverage oracle need not select the same candidate.

For state $s$, candidate set $\mathcal A_s$, and task-specific executed-outcome distance $d_s$, selection regret is
\begin{equation}
 r_s=d_s\!\left(\operatorname*{arg\,min}_{a\in\mathcal A_s}J(a)\right)
       -\min_{a\in\mathcal A_s}d_s(a).
\end{equation}
The diagnostic normalizes Cube Euclidean position distance by 0.04, Reacher maximum absolute joint-coordinate error by 0.05, and TwoRoom Euclidean position distance by 16. Its Push-T distance takes the maximum of the norm over the first four recorded position fields divided by 20 and circular angle error divided by $\pi/9$. This records the actual diagnostic implementation; its distance is not silently substituted for the closed-loop success predicate. Formal evaluation must document the state-field mapping and align the intended physical quantities.

The Phase1.5 snapshot reports mean regrets of 1.304, 1.323, and 0.232 for Push-T, Reacher, and TwoRoom. Their state-clustered 95\% intervals are [1.080, 1.540], [1.087, 1.583], and [0.219, 0.244], respectively. Different task dynamics and state definitions preclude a universal interpretation of these magnitudes. The incomplete Cube pool and control pools requiring protocol reconciliation are excluded from corresponding claims.

\paragraph{Real refinement and intervention.}
For identical starting actions, compare the changes in predicted cost and executed physical distance after PO or GF. Report valid paired counts, termination mismatches, and the rate of predicted improvement with real degradation. Matched-displacement random perturbations test whether the gradient direction contributes beyond movement magnitude. Action-order changes, local perturbations, and physical-zero controls test response to actions; zero in normalized coordinates is not assumed to be physical zero. These interventions require reliable state restoration and aligned action semantics.

\paragraph{Physical readouts.}
Readout train/validation splits are made by source trajectory, excluding diagnostic trajectories. Ridge and small-MLP probes on frozen features, together with a random-encoder control, assess state readability. Freeze the trained readout before applying it to predicted latents. Report physical quantities and units individually; aggregate errors across heterogeneous units are not a common physical scale. Three MLP seeds concern readout fitting only, not world-model training. Strong state readout without improved ranking or real refinement does not establish a coupling mechanism.

\section{Budget and Reporting Protocol}
The initial latency comparison should cover P0, P3, P0-PO, multi-step GF, and several P1/P2 CEM budgets. Synchronize device work, record warm-up and measurement counts, and isolate the single-environment decision path. Include feature encoding, candidate generation, scoring, backward passes, search updates, and action conversion. Report goal caching, precision, compilation, batch splitting, and resident memory. Environment stepping and multi-environment throughput are separate measurements. A reduction in sequential rollout depth alone is not a measured speedup or a new hardware algorithm.

\TODO{Final supplement: add resolved training histories, data exclusions, action-bound alignment, per-seed results, independent-cohort identities, completed paired refinement diagnostics, and audited timing curves.}
